% !TEX TS-program = xelatex
% !TEX encoding = UTF-8 Unicode
% !Mode:: "TeX:UTF-8"

\documentclass{resume}
\usepackage{zh_CN-Adobefonts_external} % Simplified Chinese Support using external fonts (./fonts/zh_CN-Adobe/)
%\usepackage{zh_CN-Adobefonts_internal} % Simplified Chinese Support using system fonts
\usepackage{linespacing_fix} % disable extra space before next section
\usepackage{cite}

\begin{document}
\pagenumbering{gobble} % suppress displaying page number

\name{吴思兴}

\basicInfo{
  \email{shugoforgml@gmail.com} \textperiodcentered\ 
  \phone{(+86) 180-6057-0063} \textperiodcentered\ 
  }
 
\section{\faGraduationCap  教育背景}
\datedsubsection{\textbf{闽江学院}, 福州}{2017 -- 2021}
\textit{学士}\ 软件工程

\section{\faUsers\ 实习/项目经历}
\datedsubsection{\textbf{Polyverse - 多链游戏生态系统}}{2024年7月 -- 至今}
\role{区块链后端开发}
\textit{项目简介：基于以太坊和WAX的多链游戏生态系统，融合高质量浏览器游戏体验与Web3功能，通过PTC代币促进游戏内交易和跨链资产互通}

\begin{itemize}
  \item 基于以太坊和WAX设计跨链游戏资产管理系统，使用Solidity开发ERC-20代币合约和游戏道具NFT交易功能
  \item 开发多链资产桥接合约，使用ReentrancyGuard防重入攻击，确保跨链资产安全转移
  \item 实现PTC代币经济机制，支持游戏内交易、质押奖励和治理投票功能
  \item 使用Golang构建游戏后端API服务，处理用户账户管理、游戏数据同步和链上交互逻辑
  \item 构建自动化合约部署流水线，支持多链环境下的合约管理和升级，提升部署效率
  \item 建立完善的测试体系，基于TDD开发模式编写单元测试，保障合约代码质量和资金安全
  \item 参与Python微服务向Golang的技术栈迁移，优化游戏服务性能和开发效率
\end{itemize}


\datedsubsection{\textbf{福州热游网络科技有限公司}}{2021年8月 -- 2025年1月}
\role{中台、游戏服务端} {golang python php linux mysql}

\begin{itemize}
  \item 负责《光之圣境》《女武神》《梦道2》等成熟游戏项目的账户服务、支付系统维护与功能迭代
  \item 主导《龙女仆》《魔法阵》《魔法老师》《暗之回响》等4款游戏账户服务器、支付系统从0到1开发
  \item 设计并维护游戏账户服务器，基于PHP+Nginx+MySQL+Linux架构，稳定支撑多款游戏用户认证与管理
  \item 开发完整的游戏支付系统，集成多渠道支付接口，保障交易安全性和稳定性
  \item 构建数据统计后台系统，采用Golang+PHP+Vue.js技术栈，为运营决策提供实时数据支持
  \item 开发运维工具链和自动化部署系统，使用Python、Golang、Qt C++等技术，提升运维效率
  \item 设计实现基于OAuth2的JWT统一认证鉴权体系，支持公司内部及部分外部系统的安全访问控制
\end{itemize}

\datedsubsection{\textbf{福州掌中云科技有限公司}}{2021年6月 -- 2021年8月}
\role{php后端开发}

\begin{itemize}
  \item 负责微信小程序阅读平台的后端API设计与开发，保障接口的高可用性与安全性
  \item 使用PHP进行业务逻辑开发，主要模块包括用户登录、书籍信息管理、阅读进度同步及支付接口对接
  \item 参与前端部分页面的开发与数据渲染，实现前后端高效协作
\end{itemize}


% Reference Test
%\datedsubsection{\textbf{Paper Title\cite{zaharia2012resilient}}}{May. 2015}
%An xxx optimized for xxx\cite{verma2015large}
%\begin{itemize}
%  \item main contribution
%\end{itemize}

\section{\faCogs\ IT 技能}
% increase linespacing [parsep=0.5ex]
\begin{itemize}[parsep=0.5ex]
  \item 编程语言: golang == solidity >  php == python == javascript... 
  \item 平台: linux、windows
  \item 云平台: 腾讯云、阿里云
\end{itemize}

\section{\faInfo\ 其他}
% increase linespacing [parsep=0.5ex]
\begin{itemize}[parsep=0.5ex]
  \item GitHub: https://github.com/shugoforgit
\end{itemize}

%% Reference
%\newpage
%\bibliographystyle{IEEETran}
%\bibliography{mycite}
\end{document}
